\documentclass[border=4pt]{standalone}
\usepackage{amsmath,amssymb,mathtools,xcolor,varwidth}
\definecolor{figaviolet}{HTML}{8E24AA}
\definecolor{basegreen}{HTML}{00A35A}
\definecolor{curveblue}{HTML}{1E6FE0}
\begin{document}
\begin{varwidth}{28em}
\large\bfseries
Here is the \textcolor{figaviolet}{first figure $AacE$ with its inscribed rank of parallelograms}, erected on three equal bases $AB$, $BC$, $CE$ along \textcolor{basegreen}{the baseline $AE$}. \textcolor{curveblue}{The curve $aE$} rises gently above the base, and each strip reaches up to meet it just as in Lemma II. We give this figure exactly three parallelograms because Lemma IV requires that the two figures carry an \emph{equal number} of strips in each rank. Call these three parallelograms $p_1$, $p_2$, $p_3$, taken in order from $A$ toward $E$; they are the building blocks whose ratios we will compare.
\end{varwidth}
\end{document}
